\documentclass[12pt,a4paper]{article}
\usepackage[a4paper, portrait, margin=1in]{geometry}
\usepackage{hyperref}
\usepackage{fancyhdr}
\usepackage{graphicx}
\usepackage{pdflscape}
\usepackage{svg}
\usepackage{wrapfig}
\usepackage{lastpage}
\usepackage[autostyle]{csquotes}

\MakeOuterQuote{"}
\graphicspath{ {./} }
\hypersetup{
	colorlinks=true,
	linkcolor=blue,
	filecolor=magenta,
	urlcolor=blue,
	pdftitle={Projeto AMS 2022/2023 - Entrega 1 - Grupo 38},
	pdfpagemode=FullScreen,
}

\pagestyle{fancy}
\fancyhf{}
\fancyhead[C]{Projeto AMS 2022/2023 - Entrega 1 - Grupo 38}
\fancyfoot[L]{\emph{Consultar notas na pág. \pageref*{notes}}}
\fancyfoot[C]{99209, 99256, 99311}
\fancyfoot[R]{Pág. \thepage{} de \pageref*{LastPage}}

\begin{document}

\begin{titlepage}
	\begin{center}

		\vspace*{0.5cm}
		\includegraphics[width=0.5\textwidth]{ist-logo.png}

		\vspace{1cm}
		\Huge
		\textbf{Projeto de AMS: Entrega 1}

		\vspace{0.5cm}
		\LARGE
		Grupo 38 \quad | \quad Turno L08 \quad | \quad 2022/2023

		\vfill
	\end{center}

	\large
	\begin{tabular}{c|c|c}
		\textbf{Nome}           & \textbf{Número} & \textbf{Est. Horas Trabalho} \\
		\hline
		Diogo Romão Cardoso     & 99209           & 19h                          \\
		João Fernandes Rocha    & 99256           & 19h                          \\
		Rafael Serra e Oliveira & 99311           & 19h
	\end{tabular}

	\vspace{5cm}
\end{titlepage}

\begin{landscape}
	\section*{A.1. Diagrama de Vista Geral do Negócio}
	\includesvg[inkscapelatex=false,width=1.55\textwidth]{../PICKTRU.svg}
\end{landscape}

\begin{landscape}
	\section*{A.2. Diagrama do Processo P-SET}
	\includesvg[inkscapelatex=false,width=1.53\textwidth]{../P-SET.svg}
\end{landscape}

\begin{landscape}
	\section*{A.3. Diagrama do Processo P-ON}
	\includesvg[inkscapelatex=false,width=1.55\textwidth]{../P-ON.svg}
\end{landscape}

\section*{Notas, \emph{Assumptions} e Escolhas de Desenho} \label{notes}
\normalsize
\textbf{A.1. Diagrama de Vista Geral do Negócio}
\begin{itemize}
	\item Optou-se por representar uma relação "servir" entre a Aplicação TEST e os Interfaces Aplicacionais para interação com os sistemas do cliente, visto que esta "fornece a sua funcionalidade" a esses elementos (durante a execução do processo P-ON).
\end{itemize}
\textbf{A.2. Diagrama do Processo P-SET}
\begin{itemize}
	\item Após validar com a docente, optou-se por representar explicitamente um processo "Executar Plano de Projeto", considerando-se que tal era uma unidade lógica relevante do processo (linhas 69-70 do Universo de Discurso), permitindo ainda uma notação mais abreviada relativamente ao uso do \emph{Data Object} "Plano de Projeto" nesse sub-processo.
\end{itemize}
\textbf{A.3. Diagrama do Processo P-ON}
\begin{itemize}
	\item Seguindo o aconselhado pela docente, optou-se por modelar o processo com um \emph{Message Start Event} em cada \emph{pool}, visto que o Universo de Discurso nada diz relativamente à origem do Plano de Trabalhos, nem se as unidades referidas o recebem em simultâneo.
	\item Conforme o exemplificado nas aulas teóricas, utilizaram-se \emph{Parallel Gateways} em pontos críticos de comunicação, prevenindo problemas de (falta de) sincronização entre os fluxos dos dois participantes, de modo a assegurar que nenhuma mensagem é perdida.
	\item Não sendo o Universo de Discurso explícito, e após validar com a docente, assumiu-se que não seria necessário a \emph{U-HW} aguardar que a \emph{U-SW} terminasse de analisar o plano de trabalhos antes da primeira se deslocar à loja, nem que seria necessário que a \emph{U-SW} aguardasse que a \emph{U-HW} terminasse de executar as suas tarefas na loja antes da primeira começar a preparação dos testes.
\end{itemize}

\end{document}